\begin{statementbox}
\textbf{C086\quad 双曲线焦距比与离心率取值范围（存在性判据）}

设标准双曲线
\[
H:\ \frac{x^2}{a^2}-\frac{y^2}{b^2}=1,\qquad a,b>0,
\quad F_1(-c,0),\ F_2(c,0),\ c^2=a^2+b^2,\ e=\frac ca>1.
\]
给定实数 $\lambda>0$，问：\textbf{双曲线上是否存在点 $P$，使}
\[
\boxed{PF_1=\lambda PF_2\ ?}
\]
\textbf{核心结论（充要条件）：}
\[
\boxed{\exists P\in H:\ PF_1=\lambda PF_2}
\quad\Longleftrightarrow\quad
\boxed{\lambda\ne1,\ \ 1<e\le\frac{1+\lambda}{|1-\lambda|}}.
\]
其中 $PF_1,PF_2$ 表示\textbf{线段长度}，且两者均正。

\textbf{关键边界：}
\begin{itemize}
\item $\lambda=1$：\textbf{绝无交点}，因为双曲线上的焦距差恒为 $2a>0$；
\item $1<e<(1+\lambda)/|1-\lambda|$：有\textbf{两个}关于 $x$ 轴对称的点；
\item $e=(1+\lambda)/|1-\lambda|$：只有\textbf{一个}点，为相应支的\textbf{顶点}；
\item $e>(1+\lambda)/|1-\lambda|$：\textbf{不存在}满足条件的点。
\end{itemize}
\textbf{几何重述（阿波罗尼斯圆）：}当 $\lambda\ne1$ 时，平面上
$PF_1=\lambda PF_2$ 的轨迹是\textbf{一条圆}，记为 $\mathcal A_\lambda$；
本结论等价于判别 $\mathcal A_\lambda$ 与双曲线 $H$ 是否相交。
两点相交、顶点相切、没有交点，分别对应离心率在上界之内、等于上界、超过上界。
注意 $\lambda=1$ 时轨迹退化为\textbf{焦点连线的垂直平分线} $x=0$，并非圆。

\textbf{一句话主线：}双曲线使两个焦半径的\textbf{差固定}；
比例让\textbf{和也确定}；焦点组成的三角形要求\textbf{和不小于焦距}。
\end{statementbox}

\clearpage
\begin{statementbox}
\textbf{同一结论的反向形式（固定离心率，求允许的焦距比）。}

固定 $e>1$，令 $\lambda=PF_1/PF_2$，则
\[
\boxed{\lambda\in\left[\frac{e-1}{e+1},\,1\right)
\ \cup\ \left(1,\,\frac{e+1}{e-1}\right]}.
\]
\begin{itemize}
\item 在\textbf{右支}：$1<\lambda\le(e+1)/(e-1)$，右顶点取最大值；
\item 在\textbf{左支}：$(e-1)/(e+1)\le\lambda<1$，左顶点取最小值；
\item 当 $P$ 沿任一支远离原点时，$\lambda\to1$，但不会等于 $1$。
\end{itemize}

\begin{center}\begin{tikzpicture}[font=\small,>=stealth]
\draw[gray!45,rounded corners=3pt] (0,0) rectangle (14.2,10.5);
\node at (7.1,9.85) {先认两条焦半径：长腰与短腰的差固定为 2a};

\begin{scope}[shift={(4.1,5.1)},scale=.62]
\draw[gray,->] (-5.2,0)--(5.3,0) node[right] {$x$};
\draw[gray,->] (0,-5.2)--(0,5.2) node[above] {$y$};

\draw[teal!75!black,thick,samples=90,domain=-4.8:4.8,variable=\v] plot ({2*sqrt(1+\v*\v/12)},{\v});
\draw[teal!75!black,thick,samples=90,domain=-4.8:4.8,variable=\v] plot ({-2*sqrt(1+\v*\v/12)},{\v});

\coordinate (Fone) at (-4,0);\coordinate (Ftwo) at (4,0);\coordinate (P) at (3,{sqrt(15)});
\draw[orange!85!black,very thick] (Fone)--(P);
\draw[purple!80!black,very thick] (P)--(Ftwo);
\draw[gray,dashed] (P)--(3,{-sqrt(15)});
\fill (Fone) circle(3pt) node[below=6pt] {$F_1$};
\fill (Ftwo) circle(3pt) node[below right=6pt] {$F_2$};
\fill (P) circle(3pt) node[above right=5pt] {$P$};
\fill[teal!75!black] (3,{-sqrt(15)}) circle(3pt) node[below right=5pt] {$P'$};
\node[orange!85!black] at (-.9,2.8) {$PF_1=8$};
\node[purple!80!black,anchor=west] at (4.8,2.2) {$PF_2=4$};
\node[below left] at (0,0) {$O$};
\end{scope}
\node[anchor=west,align=left] at (9.1,7.6) {图取 $a=2,c=4$\\[10pt]$e=2,\quad\lambda=2$};
\node[anchor=west,align=left] at (9.1,5.4) {差：$8-4=4=2a$\\[12pt]和：$8+4>8=2c$};
\node[anchor=west,align=left] at (9.1,3.2) {上下两点都满足\\[10pt]$PF_1=2PF_2$};
\node at (7.1,.65) {橙色始终表示 $PF_1$，紫色始终表示 $PF_2$；长度比不是向量比。};

\end{tikzpicture}
\end{center}

\textbf{顶点的精确比值：}
\[
P=(a,0):\ \frac{PF_1}{PF_2}=\frac{c+a}{c-a}=\frac{e+1}{e-1},
\qquad
P=(-a,0):\ \frac{PF_1}{PF_2}=\frac{c-a}{c+a}=\frac{e-1}{e+1}.
\]
\textbf{适用信号：}题设出现“\emph{双曲线上存在一点}”“\emph{到两个焦点距离的比}”“\emph{离心率取值范围}”。
\end{statementbox}
